  \node[draw=cblue,fill=clightblue,rounded corners=2pt,align=center,
        inner xsep=6pt,inner ysep=3.2pt,minimum width=2.75cm] (pick) at (0,-2.78)
        {copy these three $+1$ roots\\margin rises from $-2$ to $+1$};
  \coordinate (p1) at ($(pick.north west)!0.22!(pick.north east)$);
  \coordinate (p2) at ($(pick.north west)!0.50!(pick.north east)$);
  \coordinate (p3) at ($(pick.north west)!0.78!(pick.north east)$);
  \draw[-{Latex[length=1.8mm]},cblue,line width=0.75pt] (r-1.south) -- (p1);
  \draw[-{Latex[length=1.8mm]},cblue,line width=0.75pt] (r0.south) -- (p2);
  \draw[-{Latex[length=1.8mm]},cblue,line width=0.75pt] (r1.south) -- (p3);
  \node at (0,-3.60) {$-2+1+1=0$ gives a tie; a third $+1$ is needed for a strict win.};
  \node[text=cgray,font=\scriptsize] at (0,-4.02)
        {The exact sorted certificate is this same argument applied to every budget.};
\end{tikzpicture}
\caption{The Feynman version of Eq.~\eqref{eq:extrema}: a copy adds one already-observed per-root difference. Sort those differences, then add the most favorable ones first.}
\label{fig:toy}
\end{figure}

\Needspace{0.42\textheight}
\section{The Hard Part: One Winner Needs One Schedule}
Pairwise reversibility is not enough for a leaderboard claim. A target can beat rival $A$ using one copied-root set and rival $B$ using another, with no single set beating both. This is the key distinction between ``a pair can flip'' and ``this system can actually become the winner.''

\begin{figure}[H]
\centering
\begin{tikzpicture}[font=\footnotesize,>=Latex]
  \node[draw=cblue,fill=clightblue,rounded corners=2pt,minimum width=2.15cm,
        minimum height=0.52cm] (sa) at (-1.55,0) {$S_A$};
  \node[below=1.8mm of sa,draw=cgreen,fill=clightgreen,rounded corners=2pt,
        minimum width=2.15cm] (aa) {target $>$ rival $A$};
  \node[below=1.0mm of aa,draw=cred,fill=clightred,rounded corners=2pt,
        minimum width=2.15cm] (ab) {target $\le$ rival $B$};
  \node[draw=corange,fill=clightorange,rounded corners=2pt,minimum width=2.15cm,
        minimum height=0.52cm] (sb) at (1.55,0) {$S_B$};
  \node[below=1.8mm of sb,draw=cgreen,fill=clightgreen,rounded corners=2pt,
        minimum width=2.15cm] (bb) {target $>$ rival $B$};
  \node[below=1.0mm of bb,draw=cred,fill=clightred,rounded corners=2pt,
        minimum width=2.15cm] (ba) {target $\le$ rival $A$};
  \coordinate (merge) at (0,-2.35);
  \draw[gray!70,line width=0.65pt,rounded corners=2pt]
        (ab.south) -- ++(0,-3mm) -| (merge);
  \draw[gray!70,line width=0.65pt,rounded corners=2pt]
        (ba.south) -- ++(0,-3mm) -| (merge);
  \fill[gray!70] (merge) circle (1.2pt);
  \node[below=4.5mm of merge,draw=black,fill=gray!7,rounded corners=2pt,
        align=center,minimum width=5.0cm,minimum height=0.78cm] (joint)
        {one common schedule $S^*$\\must beat $A$ and $B$ simultaneously};
  \draw[-{Latex[length=1.9mm]},gray!75,line width=0.75pt] (merge) -- (joint.north);
  \node[below=1.8mm of joint,font=\bfseries] {pairwise feasible $\not\Rightarrow$ jointly attainable};
\end{tikzpicture}
\caption{Why \copycat{} needs a joint winner certificate. Separate pairwise witnesses may be mutually incompatible.}
\label{fig:joint}
\end{figure}

For binary accuracy, roots can be compressed by their outcome pattern $p$ across systems. Let $N_p$ be the number of roots with pattern $p$, $x_p$ the number copied, $C_t$ the target's original correct count, and $a_{t,p}$ the target's bit on pattern $p$. A target $t$ is a unique winner only if the \emph{same} integer vector $x$ satisfies
\begin{equation}
\sum_p x_p(a_{t,p}-a_{j,p})\ge 1-(C_t-C_j),\quad \forall j\ne t,
\label{eq:joint}
\end{equation}
with $0\le x_p\le N_p$, $x_p\in\mathbb Z$, and $\sum_p x_p\le B$. Feasible witnesses are recomputed with integer arithmetic; tiny instances are independently brute-forced.

\Needspace{0.62\textheight}
\section{Certificate Atlas on the Pinned MMLU Archive}
The most useful way to read the ten-system result is not as a leaderboard but as a map of what the allowed intervention can and cannot reach. Figure~\ref{fig:atlas} shows the original counts and the exact winner certificate for every configuration.

\begin{figure}[H]
\centering
\begin{tikzpicture}[x=0.073cm,y=0.48cm,font=\scriptsize]
  \def\xzero{488}
  \newcommand{\atlasrow}[5]{%
    \node[anchor=east,font=\scriptsize] at (-1.1,#1) {#2};
    \draw[gray!16] (0,#1) -- (46,#1);
    \pgfmathsetmacro{\len}{#3-\xzero}
    \fill[#5!72] (0,#1-0.22) rectangle (\len,#1+0.22);
    \node[anchor=west,font=\bfseries\scriptsize] at (\len+0.7,#1) {#3};
    \node[anchor=west,text=#5,font=\scriptsize,align=left,text width=1.55cm] at (51,#1) {#4};
  }
  \atlasrow{9}{Gemini 3.8 Flash}{532}{min 0}{cgreen}
  \atlasrow{8}{GPT 6.1 Sol}{530}{min 3}{cgreen}
  \atlasrow{7}{Claude Sonnet 5.5}{525}{min 8}{cgreen}
  \atlasrow{6}{DeepSeek V4.1 Flex}{511}{max $-11$}{cred}
  \atlasrow{5}{GLM 5.3 Flash Flex}{508}{max $-12$}{cred}
  \atlasrow{4}{Gemma 4 31B}{507}{max $-11$}{cred}
  \atlasrow{3}{GLM 5.3}{506}{max $-17$}{cred}
  \atlasrow{2}{Qwen 3.8 27B}{502}{max $-17$}{cred}
  \atlasrow{1}{GPT 6 Luna}{500}{max $-21$}{cred}
  \atlasrow{0}{Claude Haiku 5.5}{493}{max $-35$}{cred}
  \foreach \x/\lab in {2/490,12/500,22/510,32/520,42/530}{
    \draw[gray!28] (\x,-0.55) -- (\x,9.45);
    \node[anchor=north,text=cgray,font=\tiny] at (\x,-0.70) {\lab};
  }
  \draw[gray!42,line width=0.5pt] (48,-0.45) -- (48,9.45);
  \node[anchor=north,font=\tiny,text=cgray] at (23,-1.05) {original correct-answer count (axis begins at 488)};
  \node[anchor=west,font=\bfseries\scriptsize] at (51,9.92) {certificate};
  \node[anchor=west,font=\tiny,text=cgray] at (51,9.52) {$\le 0$ max $\Rightarrow$ excluded};
\end{tikzpicture}
\caption{Complete winner atlas for the pinned 588-root MMLU-ProX-Lite archive. Exactly three configurations are attainable unique winners under the allowed one-extra-copy-per-root intervention. GPT needs exactly 3 copies, Claude exactly 8, and the other seven receive solver-independent blocking-rival exclusions. ``Max'' is the best possible target-minus-Gemini correct-count advantage after copying every root favorable to that target; a nonpositive value excludes a unique win.}
\label{fig:atlas}
\end{figure}

The three-copy GPT witness uses root indices $\{277,507,424\}$ and yields final correct counts $533$ for GPT and $532$ for Gemini. The eight-copy Claude witness uses $\{322,380,190,424,393,245,15,564\}$ and yields Claude $533$, Gemini $532$, GPT $531$. These are not merely solver outputs: the initial deficits imply the lower bounds of three and eight, and the supplied schedules attain them.

\Needspace{0.38\textheight}
\section{Attainability Is Not Frequency}
A targeted certificate answers ``can this happen under the declared rule?'' A random control answers ``how often does it happen under one specific randomization?'' Those are different questions.

\begin{figure}[H]
\centering
\begin{tikzpicture}[font=\footnotesize,x=0.055cm,y=0.8cm]
  % MMLU stacked bar
  \node[anchor=east] at (-2,2.5) {MMLU};
  \fill[cred!80] (0,2.2) rectangle (53.3,2.8);
